\section{Information categories and relation to earlier inverses}
\label{sec:comparison}
Active boundary estimation is an essential comparison for the finite
theorem, not merely an analogy with passive billiard spectra.
Castro and Nowak~\cite{CN} study adaptive binary classification under
boundary-fragment smoothness and noise assumptions, with matching
orders for their excess-risk criterion. Their input selects a feature
point whose response is a class label. This directly available label
is not our collision bit: a solid start and a free miss have the same
output here. Neither their observation model nor their loss identifies
our finite periodic table. We use an independent physical packing and
binary-tree proof rather than importing a classification lower bound.

Locatelli, Carpentier and Kpotufe~\cite{LCK} give an adaptive strategy
for smooth decision boundaries without prior knowledge of smoothness
and noise parameters. Their membership-query setting and local
line-search/interpolation construction make the closest algorithmic
comparison. In contrast, our smoothness is known, our error is a $C^2$
geometry norm, and the boundary label must first be computed from
reciprocal pooled collision means. We claim neither a new general
bisection principle nor an advance on adaptation to unknown smoothness.
The independent reduction in Proposition~\ref{prop:query} is what
permits this established active-estimation strategy in the physical
collision experiment.

The earlier scalar inverse reconstructed a uniformly accurate occupation
on a full two-dimensional grid. Its constructive upper bound was
$C\nu^{-3}\log(C/(\nu\delta))$ attempts on the same bounded
$C^{6,\beta}$ class, using $\sigma\asymp\nu^{3/2}$ and compensated
support smoothing. That theorem remains in Supplement R. The present
procedure uses $O(h^{-1})$ angular lines, a fixed-depth dependency
stencil for each queried point, and the full $C^{6,\beta}$ regularity.
It has a different adaptive command design, not a sharper analysis of
the old uniform design. We have not proved a matching lower bound
for that old design or for uniform spatial preparations.

The resource improvement is accompanied by the explicit localization
and calibration cost in \eqref{eq:precision}. Exact data remain a pair
of functionals on controlled spatial inputs. The lower bound concerns
only their finite binary realizations; it does not assert a low-information
passive invariant. The angle/time/position tolerances shrink with
localization and are not learned from the collision outcomes. Apparatus
motion, precision in representing input locations, and arithmetic
cost are distinct resources.

For completeness the exact v31 manuscript is supplied unchanged as
Supplement R. It retains the arbitrary centered-law and general
nontrivial bounded-law inverses, the different-zero-set comparison,
the full fixed-aperture proof, rational interval enclosures, the
uniform-grid construction, and all earlier moment, marked-law and
count-fiber results in their original nested volumes. The more general
randomized law is an exact functional theorem; the finite local query
proved here uses the four-atom compass law. No continuous transition
law is discretized without an error budget. Rational intervals remain
conditional on forcing and survival bounds, not certificates of the
apparatus or physical prior.

The finite global conclusion also retains the known positive patch
margin and the bounded periodic presentation. Exact rational period
relations are distinct from approximate Euclidean generators. Removing
the known margin gives only pointwise eventual recognition, as described
in Section~\ref{sec:periods}. The result does not settle rigidity from
uniformly prepared count germs, exact analytic count germs, marked
length spectra or passive trajectories. These boundaries delimit the
information theorem without changing its topic or weakening any retained
mathematical assertion.
